\section{Reproducibility and review}
The repository supplies dependency-free Python experiments, unit
tests, LaTeX sources, references, a claim register, and CI configuration.
The experiment constructs the truth table of \(g_m\), enumerates every
prefix and suffix in each selected order, and counts distinct residual
truth tables. The published range is \(1\leq m\leq7\).
All arithmetic is exact; there are no random seeds or statistical
estimates.

Tests compare the incidence predicate with an independent exhaustive
unit-clause SAT verifier, check the grouped and interleaved width
predictions on small inputs, and test a separate parity example.
The manuscript tables are regenerated and checked for drift.
These checks validate finite artifacts, not the general mathematical
proofs.

For graph selectors, a second exact experiment counts residuals for every
read-variable subset. Their number depends on that subset, not on its internal
order. A minimax dynamic program over subsets finds an optimal variable order,
including initial and terminal layers. It evaluates matching-3, cycle-6,
grid-2x3, cube-3, and grid-3x3, with at most nine input bits.
Tests compare the optimizer with independent exhaustive permutation enumeration
on selected functions of at most five bits. The graph predicate is also checked
against a witness-enumerating SAT verifier for every simple graph on at most
four vertices. These finite graphs include an odd grid; the data are not a
proof of Proposition~\ref{prop:grid-orders}, which assumes an even side length.

The command \texttt{make check} runs code and research-metadata checks;
\texttt{make pdf} builds the paper with \texttt{latexmk};
\texttt{make formal} verifies Lemma~\ref{lem:residual} and
Proposition~\ref{prop:grouped} in pinned Lean~4.35.0;
\texttt{make verify} performs all these checks and checks the LaTeX log.
The README lists the toolchain.
The formal project uses the bundled standard library without Mathlib.
It proves the residual-width lower bound and one residual diagram attaining
all cuts, including the initial and final layers. An axiom audit permits
only Lean's standard logical axioms for these theorems; it rejects proof
placeholders and custom axioms. The repository's formal guide matches the
finite reachable-state and fixed-length conventions to the manuscript.
The grouped-disjointness formalization also proves the exact midpoint count
\(2^m\) for every \(m\geq0\), including the empty-input case, and obtains
the minimum-width characterization from the residual theorem. These
formalizations cover L02 and P03, not the other claims.
Independent review must inspect the evaluator model, uniformity
quantifiers, finite exceptions, and representation restrictions.
Journal or preprint submission remains an author-controlled future step.

\paragraph{Authorship and assistance.}
The research question was developed by Matías Donoso Quiroga.
OpenAI ChatGPT/Codex assisted in preparing this initial exposition,
code, and repository structure. Independent human mathematical
review is pending. AI systems are not listed as authors.
No institutional affiliation, publication venue, or endorsement
is asserted.
